Recent work in deep learning has demonstrated a strong interest in producing embeddings
on the unit hypersphere as opposed to an unconstrained latent space. Downstream tasks
such as retrieval, clustering, and k-NN classification rely on angular relationships
between embeddings, and outside of uncertainty \citep{levi2025doubleellipsoidgeometryclip},
the $\ell^2$ norm of an embedding tends to convey little information. Variational autoencoders (VAEs) serve as a natural testbed for investigating how hyperspherical
embeddings differ from, or improve on, Euclidean embeddings in the context of deep generative models.

The generative model for a VAE is defined as follows:
\[
    \begin{cases}
        \mathbf{z} &\sim p(\mathbf{z}) \\
        \mathbf{x} \mid \mathbf{z} &\sim p_{\boldsymbol{\theta}}(\mathbf{x} \mid \mathbf{z})
    \end{cases}
\]
and the variational approximation to the posterior is
\[
    \mathbf{z} \mid \mathbf{x} \sim q_{\boldsymbol{\phi}}(\mathbf{z} \mid \mathbf{x})
\]
where $\mathbf{x}$ denotes the observed data and $\mathbf{z}$ denotes the latent variables. The subscripts
$\boldsymbol{\theta}, \boldsymbol{\phi}$ denote parametrization by a deep neural network. In a Euclidean latent
space, the usual choice for prior $p(\mathbf{z})$ is $\mathcal{N}(0, I)$. Ideally, one would want to choose a
prior that is uniform over the support of $\mathbf{z}$, thus motivating the need for a compact latent space. On the
hypersphere, $\text{Uniform}(\mathbb{S}^{d - 1})$ is then a natural choice of prior. As for extending the posterior
to the hypersphere, the vMF distribution is a natural adaptation of the normal distribution. Parametrized by a mean
direction $\boldsymbol{\mu}$ and concentration parameter $\kappa$, it can be considered a ``location-scale'' family
in a similar sense to the normal distribution. \cite{davidson2022hypersphericalvariationalautoencoders} have shown that
this hyperspherical generative model setup can match, or in some cases improve on, Gaussian performance on simple tasks
in low dimensions.

The vMF distribution does not come without its own problems, though. Its density involves the modified Bessel function
which becomes unstable for large inputs, and it may only be sampled through rejection sampling which incurs a large
speed hit for large values of $\kappa$. To mitigate these issues, \cite{decao2020powersphericaldistribution} introduced
the power-spherical distribution, which is both numerically stable and can be sampled through a deterministic transformation
of a beta-distributed sample. Even so, neither distribution can admit more than one mode on the unit hypersphere, which
limits their expressivity. To that end, this paper introduces the SphericalTNBBeta distribution which exhibits the same
benefits of the power-spherical distribution in the context of a VAE, but may also express two, or sometimes three modes.
In evaluating hyperspherical VAEs on a series of downstream tasks, we find that the SphericalTNBBeta distribution is particularly
well-positioned to represent featureless graphs, and performance improvement tends to correlate with the number of overlapping
communities.

